\documentclass[10pt]{article}
\usepackage[margin=1in]{geometry}
\usepackage{amsmath,amssymb,booktabs,graphicx,microtype,hyperref}
\usepackage[ruled,vlined,linesnumbered]{algorithm2e}
\usepackage{placeins}
\usepackage{needspace}
\usepackage{tabularx}
\usepackage{siunitx}
\usepackage{xcolor}
\graphicspath{{figures/}}

\title{ORBIT: Function-Space Optimization for Rotary Query--Key Interactions}
\author{Alexandru Pop\\[-0.1em]\small Independent Researcher\\[-0.1em]\small \href{mailto:alexandrupp55@gmail.com}{alexandrupp55@gmail.com}}
\date{}

\hypersetup{
  pdftitle={ORBIT: Function-Space Optimization for Rotary Query--Key Interactions},
  pdfauthor={Alexandru Pop},
  pdfsubject={RoPE-aware query--key optimization for language-model training},
  pdfkeywords={ORBIT, Muon, RoPE, optimization, attention, language models}
}

\IfFileExists{generated/macros.tex}{\input{generated/macros.tex}}{}

\begin{document}
\maketitle

\begin{abstract}
Matrix optimizers such as Muon exploit the spectral structure of parameter matrices while
treating query and key projections as ordinary matrix parameters. ORBIT introduces a
function-space extension of Muon for rotary attention: frequency-local $2\times2$ metrics
are constructed from query/key covariances transported through RoPE relative-position
rotations, applied to the candidate Q/K updates, and followed by joint Frobenius-norm
restoration. On RoPE-modified GPT-2-style decoders trained on FineWeb-Edu, the evaluation
separates optimizer mechanism from hyperparameter recipe. Under an identical configuration
selected independently by both methods from the same ten-candidate grid, ORBIT lowers
validation loss on \MatchedWins{} held-out seeds, with a paired mean difference of
\MatchedDelta{} nats (95\% CI [\MatchedCILow{}, \MatchedCIHigh{}]). Mean training
runtime increases by \MatchedRuntimeOverhead{}, with essentially unchanged peak allocated
memory. Cross-configuration experiments show an ORBIT improvement under both frozen
recipes while revealing substantial sensitivity to recipe choice. Expanded ablations
identify contributions from functional Q/K conditioning and RoPE-aware rotation.
Longer-horizon and 355M experiments further show transfer beyond the primary setting
against Muon and NorMuon in the tested seeds. ORBIT changes no inference-time computation.
\end{abstract}

\section{Introduction}
Matrix optimizers have become a practical alternative to coordinate-wise adaptive methods
for Transformer training. Muon, for example, applies momentum followed by approximate
matrix orthogonalization to hidden two-dimensional parameters, and its scaled variants have
shown that spectral update geometry can remain effective at large language-model
scale~\cite{jordan2024muon,liu2025muon}. Yet this treatment is largely parameter-centric:
two matrices with the same shape receive the same class of update even when they participate
in qualitatively different computations.

Self-attention provides a particularly sharp example. Query and key projections jointly
determine pre-softmax attention scores, and RoPE rotates two-dimensional Q/K coordinate
pairs so that their interaction depends on relative position~\cite{su2021roformer}. For a
frequency pair $f$, the relevant score contribution is
$q_{i,f}^{\top}R_f(i-j)k_{j,f}$. Consequently, an equal-sized perturbation of $W_Q$ or
$W_K$ can induce very different functional movement depending on the opposite circuit's
statistics and on the RoPE rotation associated with token displacement. Standard Muon does
not model this interaction.

ORBIT incorporates that functional geometry without replacing Muon's efficient spectral
update or changing inference. It first forms the ordinary
Muon-style candidate Q/K directions, then preconditions each RoPE frequency pair with a
compact metric derived from the covariance of the opposite circuit after transport through a
fixed distribution of relative-position rotations. A single joint Frobenius restoration
prevents this conditioning from winning merely by increasing the total Q/K step norm. The
result is an optimizer-only modification whose unique operations are local $2\times2$
matrix functions.

Optimizer comparisons are unusually sensitive to tuning policy, making experimental design
part of the contribution. The evaluation first crosses the frozen Muon and ORBIT recipes,
then presents both methods with the same candidate set and evaluates their independently
selected winner on unseen paired seeds. This mechanism-isolated comparison anchors the
subsequent ablations and longer-horizon/model-scale transfer experiments.

\Needspace{9\baselineskip}
\paragraph{Contributions.}
The main contributions are:
\textbf{(1) RoPE-conditioned Q/K geometry}, a frequency-local functional metric derived
from the relative-position attention interaction;
\textbf{(2) matched optimizer evaluation}, separating the ORBIT update from recipe
selection across unseen seeds;
\textbf{(3) mechanism ablations}, isolating functional conditioning, RoPE-aware rotation,
and full frequency-pair coupling; and
\textbf{(4) transfer evidence}, extending the evaluation to a longer training horizon and
a 355M model with measured runtime and memory cost.

\section{Related Work}
\paragraph{Matrix and spectral optimization.}
Muon replaces a raw matrix-gradient step with a momentum-smoothed, approximately
orthogonalized direction~\cite{jordan2024muon}; later work showed that weight decay and
update-scale calibration make this family viable at substantially larger language-model
scale~\cite{liu2025muon}. PRISM adds anisotropic spectral shaping to Muon-like first-order
updates using low-rank curvature information~\cite{yang2026prism}. These methods motivate
ORBIT's spectral base update, but they do not specialize the metric to the functional role
of rotary Q/K interactions.

\paragraph{Optimization under architectural geometry.}
Recent work increasingly argues that optimizers should respect invariances or induced
weight-space geometry. \emph{Riemannian Gradient Descent for Low-Rank
Architectures}~\cite{knight2026} treats QK/VO products as low-rank optimization objects.
\emph{Dead-Direction Conditioners}~\cite{shirodkar2026ddc} develops gauge-equivariant
conditioners for architectural symmetries, including a per-head attention rotation matched
to RoPE. \emph{The Loss Does Not See the Basis, but Adam Does}~\cite{singh2026lossbasis}
shows that optimizer behavior can depend on arbitrary factor bases even when the represented
function is unchanged, with Q/K products serving as attention invariants. In
parameter-efficient adaptation, LoRA-TSD and ISO-LoRA optimize the induced tangent-space
weight update rather than treating low-rank factors independently
~\cite{andriianov2026loratsd,zhu2026isolora}. ORBIT is complementary: it constructs a local
training metric directly from the RoPE-conditioned Q/K bilinear function.

\paragraph{Query--key and RoPE structure.}
\emph{Faster Query-Key Learning Sharpens Attention}~\cite{vashisht2026fasterqk} shows
that the relative learning speed of QK and OV circuits changes attention concentration,
reinforcing the distinct functional role of QK optimization. \emph{Spectral Query-Key
Product Weight Steering}~\cite{tiwari2026qksteering} directly edits singular modes of the
query--key product as a training-free intervention. RoFormer introduced the rotation
structure that turns absolute Q/K rotations into relative-position interactions
~\cite{su2021roformer}; LeRoPE makes those frequencies learnable
~\cite{karypis2026lerope}; and \emph{RoPE Attention Is an Exact Forward-Pass Gradient Step
with Softmax Intact}~\cite{huang2026ropegradient} derives an exact gradient-step
representation of the RoPE-softmax forward pass. These works analyze or modify attention
geometry, but do not use the relative-position rotation itself to precondition Q/K
parameter updates during pretraining.

\paragraph{Positioning.}
The literature audit did not identify prior work introducing a training optimizer whose Q/K
preconditioner is derived from the RoPE-conditioned relative-position bilinear interaction
and implemented with per-frequency $2\times2$ statistics. ORBIT's novelty is therefore
specific to this training-time metric and its mechanism-isolated evaluation, rather than to
spectral descent, QK-product optimization, architecture-aware optimization, RoPE geometry,
or gauge-equivariant conditioning.

\section{Method}

\subsection{Muon matrix update}
ORBIT retains Muon as the base update for matrix-valued parameters. For a matrix parameter
$W_t$ with gradient $G_t$ and momentum state $m_t$, the implementation forms
\begin{equation}
m_t=\mu m_{t-1}+G_t,\qquad
U_t=\alpha(W_t)\,\operatorname{NS}_5\!\left(G_t+\mu m_t\right),
\label{eq:muon-candidate}
\end{equation}
where $\mu=0.95$, $\operatorname{NS}_5$ denotes five iterations of the quintic
Newton--Schulz polar map, and
$\alpha(W)=\sqrt{\max(1,n_{\mathrm{row}}/n_{\mathrm{col}})}$ is the standard
aspect-ratio scaling used by the Muon implementation~\cite{jordan2024muon}. ORBIT changes
this candidate only for query and key projections. Other two-dimensional hidden operators
use the ordinary Muon update, while embeddings, the tied language-model head, biases, and
normalization vectors use the auxiliary AdamW path described in
Section~\ref{sec:training-setup}.

\subsection{RoPE-conditioned query--key geometry}
For a single attention head and RoPE frequency pair $f$, the pre-softmax score contribution
between positions $i$ and $j$ is
\begin{equation}
  s_{ij,f}=q_{i,f}^{\top}R_f(\Delta)k_{j,f},\qquad \Delta=i-j.
  \label{eq:rope-score}
\end{equation}
A perturbation of the query therefore interacts with the rotated key statistics, while a
key perturbation interacts with the corresponding rotated query statistics. Let
$C_{Q,f}$ and $C_{K,f}$ denote exponential-moving-average $2\times2$ covariances of the
unrotated query and key coordinates. ORBIT defines
\begin{align}
M_{Q,f} &= \mathbb{E}_{\Delta}\!\left[R_f(\Delta)C_{K,f}R_f(\Delta)^\top\right],\\
M_{K,f} &= \mathbb{E}_{\Delta}\!\left[R_f(\Delta)^\top C_{Q,f}R_f(\Delta)\right].
\label{eq:orbit-metrics}
\end{align}
The expectation is evaluated on the displacement grid
$\mathcal{D}=\{1,2,4,8,16,32,64,128\}$. Covariances are updated during training with EMA
coefficient $\beta=0.95$ after the first observation. These statistics are optimizer-only
buffers and are absent from the inference graph.

\subsection{ORBIT preconditioning}
Let $U_Q$ and $U_K$ be the Muon candidates from Equation~\ref{eq:muon-candidate}. For each
head and RoPE frequency pair,
\begin{align}
\widehat U_{Q,f} &= M_{Q,f}^{-1/2}U_{Q,f},\\
\widehat U_{K,f} &= M_{K,f}^{-1/2}U_{K,f}.
\end{align}
The metrics are regularized by $10^{-5}I$, their effective condition number is capped at
100, and the inverse square root is evaluated analytically for each symmetric positive
$2\times2$ matrix.

A single joint restoration factor then preserves the total Q/K update budget:
\begin{equation}
\rho =
\sqrt{
\frac{\|U_Q\|_F^2+\|U_K\|_F^2}
     {\|\widehat U_Q\|_F^2+\|\widehat U_K\|_F^2}
},
\qquad
\widetilde U_Q=\rho\widehat U_Q,\quad
\widetilde U_K=\rho\widehat U_K.
\label{eq:joint-restore}
\end{equation}
The functional metric therefore changes the direction of the joint Q/K update without
increasing its aggregate Frobenius norm.

Figure~\ref{fig:orbit-overview} summarizes the data flow of the optimizer before the
step-by-step specification in Algorithm~\ref{alg:orbit}.

\IfFileExists{figures/orbit_overview.pdf}{%
\begin{figure}[t]
\centering
\includegraphics[width=0.98\linewidth]{orbit_overview.pdf}
\caption{ORBIT update pathway. Muon candidate directions are reshaped using Q/K covariance
statistics transported through RoPE relative-position rotations, producing frequency-local
$2\times2$ metrics. Analytic inverse-square-root preconditioning is followed by joint
Frobenius-norm restoration before the parameter update. Other parameter classes retain
their standard Muon or auxiliary AdamW path.}
\label{fig:orbit-overview}
\end{figure}}{}

Algorithm~\ref{alg:orbit} gives the corresponding update for one attention layer. Standard
Muon and auxiliary AdamW updates proceed unchanged for the remaining parameter classes.

\begin{algorithm}[t]
\caption{ORBIT update for a RoPE query--key projection pair}
\label{alg:orbit}
\DontPrintSemicolon
\KwIn{$W_Q,W_K$; gradients $G_Q,G_K$; momentum $m_Q,m_K$; covariances $C_Q,C_K$;
learning rate $\eta$; weight decay $\lambda$; momentum $\mu$; EMA $\beta$;
displacements $\mathcal{D}$}
\textbf{Forward statistics:} update $C_Q,C_K$ from current unrotated Q/K activations
with EMA coefficient $\beta$\;
$m_Q \leftarrow \mu m_Q+G_Q$;\quad $m_K \leftarrow \mu m_K+G_K$\;
$U_Q \leftarrow \operatorname{MuonCandidate}(G_Q+\mu m_Q)$;\quad
$U_K \leftarrow \operatorname{MuonCandidate}(G_K+\mu m_K)$\;
\ForEach{attention head $h$ and RoPE frequency pair $f$}{
  $M_{Q,f}\leftarrow |\mathcal{D}|^{-1}\sum_{\Delta\in\mathcal{D}}
  R_f(\Delta)C_{K,f}R_f(\Delta)^\top + 10^{-5}I$\;
  $M_{K,f}\leftarrow |\mathcal{D}|^{-1}\sum_{\Delta\in\mathcal{D}}
  R_f(\Delta)^\top C_{Q,f}R_f(\Delta) + 10^{-5}I$\;
  $T_{Q,f}\leftarrow \operatorname{InvSqrt}_{2\times2}(M_{Q,f})$;\quad
  $T_{K,f}\leftarrow \operatorname{InvSqrt}_{2\times2}(M_{K,f})$\;
  $\widehat U_{Q,f}\leftarrow T_{Q,f}U_{Q,f}$;\quad
  $\widehat U_{K,f}\leftarrow T_{K,f}U_{K,f}$\;
}
$\rho\leftarrow
\sqrt{(\|U_Q\|_F^2+\|U_K\|_F^2)/
(\|\widehat U_Q\|_F^2+\|\widehat U_K\|_F^2)}$\;
$\widetilde U_Q\leftarrow\rho\widehat U_Q$;\quad
$\widetilde U_K\leftarrow\rho\widehat U_K$\;
$W_Q\leftarrow(1-\eta\lambda)W_Q-\eta\widetilde U_Q$;\quad
$W_K\leftarrow(1-\eta\lambda)W_K-\eta\widetilde U_K$\;
\end{algorithm}

\subsection{Cost and ablation variants}
The additional state consists of two $2\times2$ covariance matrices for every attention
head and RoPE frequency pair. The unique per-step matrix operations are local $2\times2$
transforms rather than dense model-width preconditioners, leaving inference unchanged.

Three controls isolate the update components. \textbf{Identity} keeps the ORBIT training
path and statistics collection but disables functional preconditioning. \textbf{No-RoPE}
uses Q/K covariance metrics without relative-position rotation. \textbf{Diagonal} retains
RoPE-aware per-frequency scaling while removing off-diagonal coupling within each
$2\times2$ pair.

\section{Experimental Setup}
\subsection{Models and data}
All reported comparisons use the same decoder implementation within a given size. The
124M model has 12 layers, 12 attention heads, width 768, and approximately 124M parameters.
The 355M model has 24 layers, 16 heads, width 1024, and uses gradient checkpointing. Both
use context length 512, GPT-2 token embeddings with tied input/output weights, pre-norm
Transformer blocks, GELU MLPs with $4d$ hidden width, and zero dropout. Learned absolute
position embeddings are replaced by RoPE with base 10,000. Attention uses separate
$q_{\mathrm{proj}}$, $k_{\mathrm{proj}}$, and $v_{\mathrm{proj}}$ matrices.

The corpus is FineWeb-Edu. Text is tokenized with the GPT-2
tokenizer. A fixed stream of 12.4M tokens is shared across runs: the first 400,000 tokens
form the validation pool and the next 12,000,000 tokens form the training pool. Every training step
samples random contiguous 512-token windows from that pool with a seeded generator;
evaluation uses a fixed, separate generator and averages 20 validation batches. Thus every
optimizer in a paired setting sees the same validation windows.

\subsection{Training setup}
\label{sec:training-setup}
All optimizers use the same warmup--cosine multiplier. The first 10\% of steps are warmup;
the remaining schedule follows a cosine decay ending at $0.1$ times the peak learning rate.
The same multiplier is applied to every learning-rate field, including the auxiliary AdamW
path. Global gradient norm is clipped at 1.0. Training uses CUDA autocast in float16 with
gradient scaling. The 124M model uses batch size 8 (4,096 sampled tokens per optimizer
step); the 355M model uses batch size 4 (2,048 sampled tokens per step). The latter
difference is disclosed because it prevents interpreting the 355M result as a pure
model-size scaling experiment.

\subsection{Baselines and hyperparameter selection}
The broad campaign includes AdamW, Muon, NorMuon, AdaMuon, ASTRO, and ORBIT. ASTRO is a
pre-existing frozen baseline selected during the preceding ASTRO development study before
any ORBIT result was observed; it is not retuned during the ORBIT campaign.

Every optimizer receives the same number of tuned dimensions in the broad campaign.
Muon-family methods search learning rate in $[2\times10^{-3},10^{-1}]$, weight decay in
$[10^{-3},3\times10^{-1}]$, and auxiliary AdamW multiplier in $[0.02,1.5]$, with
log-uniform deterministic sampling. AdamW tunes learning rate, weight decay, and $\beta_2$.
AdaMuon uses its published natural learning-rate range but the same weight-decay and
auxiliary-rate policy. Five discovery trials are used per optimizer in the broad campaign.

The primary Muon--ORBIT comparison uses a stricter matched policy: ten deterministic
900-step candidate configurations are generated once from the Muon search space and
presented identically to both methods. Both methods independently selected the same candidate from the shared grid before any
held-out matched-confirmation seed was evaluated. That selected setting is reused in
the matched confirmation and expanded ORBIT ablation.

\subsection{Evaluation protocol}
The evaluation is organized into complementary study blocks rather than a single benchmark:
broad optimizer context, matched mechanism isolation, component ablations, and transfer
across training horizon and model size. Table~\ref{tab:protocol} summarizes the design.

\begin{table}[t]
\centering
\caption{Experimental design. The primary estimate is the matched Muon--ORBIT comparison;
the remaining experiments test recipe interaction, mechanism attribution, or transfer.}
\label{tab:protocol}
\small
\setlength{\tabcolsep}{4pt}
\renewcommand{\arraystretch}{1.12}
\begin{tabularx}{\textwidth}{@{}
>{\raggedright\arraybackslash}p{2.05cm}
>{\raggedright\arraybackslash}X
>{\centering\arraybackslash}p{0.95cm}
>{\centering\arraybackslash}p{0.85cm}
>{\centering\arraybackslash}p{1.35cm}
>{\raggedright\arraybackslash}p{3.1cm}
@{}}
\toprule
Experiment & Scientific purpose & Model & Steps & Replication & Compared methods \\
\midrule
Broad benchmark &
Frozen-recipe optimizer comparison &
124M & 900 & $n=5$ &
AdamW, AdaMuon, Muon, NorMuon, ASTRO, ORBIT \\

Crossed recipes &
Separate update-rule and recipe effects &
124M & 900 & $n=5$/cell &
Muon, ORBIT \\

\textbf{Matched confirmation} &
\textbf{Mechanism-isolated primary comparison} &
\textbf{124M} & \textbf{900} & \textbf{$n=10$ paired} &
\textbf{Muon, ORBIT} \\

Ablation &
Attribute the functional, RoPE, and pair-coupling components &
124M & 900 & $n=10$/control &
ORBIT, identity, no-RoPE, diagonal \\

Long-horizon transfer &
Test persistence under longer optimization &
124M & 2700 & $n=2$ &
Muon, NorMuon, ASTRO, ORBIT \\

Scale transfer &
Test transfer to a larger decoder &
355M & 900 & $n=2$ &
Muon, NorMuon, ASTRO, ORBIT \\
\bottomrule
\end{tabularx}

\vspace{0.35em}
\footnotesize
Broad hyperparameter selection used five discovery candidates per optimizer. The matched
comparison used ten shared candidate configurations selected before the held-out runs. The
final ablation followed a three-run pilot and used ten paired runs per control.
\end{table}

At 124M, 300, 900, and 2700 optimizer steps correspond to approximately 1.23M, 3.69M,
and 11.06M sampled training tokens for the discovery, primary, and long-horizon regimes.
The 355M transfer uses 1.84M sampled training tokens at 900 steps because of its smaller
batch size.

\paragraph{Statistical analysis.}
The primary estimand is the paired per-run validation-loss difference, ORBIT minus the
comparison method. Reported statistics include the paired mean, sample standard deviation,
paired-run win count, and a two-sided 95\% Student-$t$ confidence interval for the mean
paired difference. The matched Muon--ORBIT comparison is the primary confirmatory contrast;
ablations provide mechanism attribution, while the long-horizon and 355M experiments assess
transfer. Wall-clock time and peak GPU-memory allocation are reported alongside loss.

All final comparisons use the same frozen implementation and a consistent run-record format.
Optimizer settings, random seeds, environment metadata, validation metrics, runtime,
memory use, and training trajectories are retained with each run. Results are merged only
after automated consistency checks verify the expected experiment structure, matched
configuration, finite metrics, and a common implementation/software state. Exact
environment details and run metadata are retained in the accompanying reproducibility
artifacts. The optimizer implementation, frozen run records, and reproduction scripts are
publicly available at \href{https://github.com/pop123-ux/ORBIT}{github.com/pop123-ux/ORBIT}.


\section{Results}

\subsection{Matched hyperparameters isolate a reproducible ORBIT effect}
The primary comparison removes optimizer-specific recipe selection as a confound. Muon and
ORBIT were presented with the same ten deterministic 900-step candidate configurations and
independently selected the same frozen candidate before any held-out confirmation seed was
evaluated. Under that shared configuration, Muon reached mean validation loss
\MatchedMuonLoss{} and ORBIT reached \MatchedOrbitLoss{} across ten new paired seeds
(Table~\ref{tab:matched}).

The paired ORBIT-minus-Muon difference was \MatchedDelta{} nats, with a 95\% confidence
interval of [\MatchedCILow{}, \MatchedCIHigh{}].
\textbf{ORBIT obtained the lower validation loss on \MatchedWins{} held-out runs}
(Figure~\ref{fig:matched}). This matched estimate directly isolates the optimizer contribution from independent recipe
selection. The improvement therefore persists when both methods share the same
hyperparameters and is not explained solely by a more favorable ORBIT recipe.

Mean wall-clock time increased by \MatchedRuntimeOverhead{} relative to Muon, while
recorded peak CUDA allocation changed by only \MatchedMemoryOverhead{}. ORBIT therefore
achieves the matched loss improvement with essentially unchanged peak allocated memory and
a measurable training-time overhead.

\IfFileExists{generated/table_matched.tex}{\input{generated/table_matched.tex}}{}

\IfFileExists{figures/matched_effect.pdf}{%
\begin{figure}[t]
\centering
\includegraphics[width=0.78\linewidth]{matched_effect.pdf}
\caption{Primary matched-hyperparameter result. Circles show paired held-out run
differences; the diamond shows the paired mean with its 95\% confidence interval. Negative
values favor ORBIT.}
\label{fig:matched}
\end{figure}}{}

\subsection{Hyperparameter recipe accounts for a substantial part of the broad gap}
The optimizer-by-configuration cross-over makes the tuning confound explicit. When both
methods use the Muon-selected recipe, ORBIT improves on Muon by
\XConfigMechanismMuon{} nats and wins \XConfigMechanismMuonWins{} paired seeds. When both
use the ORBIT-selected recipe, the corresponding mechanism difference is
\XConfigMechanismOrbit{} nats, again with ORBIT winning
\XConfigMechanismOrbitWins{} seeds. Thus the ORBIT update improves the objective under both
recipes, but the size of the gain depends strongly on the surrounding hyperparameters.

Recipe choice itself has a larger effect than the matched mechanism contrast. Moving Muon
from its own discovery recipe to the ORBIT-selected recipe changes validation loss by
\XConfigRecipeMuon{} nats, while the analogous change for ORBIT is
\XConfigRecipeOrbit{} nats. The cross-over shows that the larger independently tuned ORBIT--Muon gap combines both
update-rule and recipe effects. \textbf{Across both frozen recipes, ORBIT retains the lower
validation loss on every paired run} (Figure~\ref{fig:mechanism}a).

\subsection{Ablations identify functional Q/K conditioning and RoPE-aware rotation}
The expanded ten-seed ablation tests which pieces of ORBIT are responsible for the matched
effect. The \textbf{identity control} removes the functional preconditioner while preserving the
same matched recipe and ORBIT update path. Full ORBIT improves on this control by
\AblIdentityDelta{} nats, with 95\% CI
[\AblIdentityCILow{}, \AblIdentityCIHigh{}], and wins \AblIdentityWins{} seeds. The
magnitude is close to the independent matched Muon comparison, providing a second
mechanism-focused estimate on a different seed set.

The \textbf{no-RoPE control}, which retains generic Q/K functional adaptation while
removing the relative-position rotations, weakens the optimizer further: full ORBIT improves on the no-RoPE control by
\AblNoRoPEDelta{} nats, with 95\% CI
[\AblNoRoPECILow{}, \AblNoRoPECIHigh{}], and wins \AblNoRoPEWins{} seeds. This comparison
supports the claim that relative-position-aware conditioning contributes beyond generic
Q/K covariance adaptation in the tested architecture.

Full ORBIT also improves on the \textbf{diagonal control} by \AblDiagDelta{} nats, with 95\% CI
[\AblDiagCILow{}, \AblDiagCIHigh{}], and wins \AblDiagWins{} seeds. This smaller contrast
indicates an additional contribution from the full frequency-pair metric beyond diagonal
per-frequency scaling.

\IfFileExists{figures/mechanism_summary.pdf}{%
\begin{figure}[!htbp]
\centering
\includegraphics[width=0.98\linewidth]{mechanism_summary.pdf}
\caption{Mechanism isolation. \textbf{(a)} Crossed optimizer--recipe experiment: ORBIT
retains lower loss under both frozen recipes. \textbf{(b)} Paired mechanism ablations with
95\% confidence intervals; negative values favor full ORBIT.}
\label{fig:mechanism}
\end{figure}}{}

\subsection{ORBIT transfers across training horizon and model size}
The transfer experiments preserve each optimizer's independently frozen recipe and evaluate
whether the observed advantage extends beyond the primary 124M/900-step setting.

At 124M and 2700 steps, \textbf{ORBIT remains below every comparison recipe in both tested
runs}.
Its mean paired differences are \HorizonVsMuon{} nats versus Muon,
\HorizonVsNorMuon{} versus NorMuon, and \HorizonVsAstro{} versus ASTRO. The two
ORBIT--ASTRO seed differences are closely aligned in sign and magnitude. The two paired seeds show that the advantage persists over a substantially longer training
horizon in this transfer setting.

At 355M and 900 steps, \textbf{ORBIT again outperforms the Muon and NorMuon recipes in
both tested runs}, with mean
paired differences of \ScaleVsMuon{} and \ScaleVsNorMuon{} nats respectively. The ASTRO
comparison is mixed across the two seeds: the mean ORBIT-minus-ASTRO difference is
\ScaleVsAstro{} nats, with one seed favoring each method. The larger model uses a smaller
batch size and hence a different sampled-token budget per optimizer step; the experiment
therefore demonstrates transfer to 355M rather than a controlled scaling law.

\IfFileExists{figures/transfer_summary.pdf}{%
\begin{figure}[!htbp]
\centering
\includegraphics[width=0.96\linewidth]{transfer_summary.pdf}
\caption{Transfer beyond the primary setting. \textbf{(a)} 124M model at 2700 steps.
\textbf{(b)} 355M model at 900 steps. Points show individual runs; horizontal segments show
the method mean.}
\label{fig:transfer}
\end{figure}}{}

\subsection{Broader optimizer context}
The independently tuned 124M confirmation provides a practical comparison across the wider
optimizer set. ASTRO reaches mean validation loss \BroadAstroLoss{}, ORBIT
\BroadOrbitLoss{}, NorMuon \BroadNorMuonLoss{}, and Muon \BroadMuonLoss{}. ORBIT and ASTRO are numerically close under their independently selected recipes, while
\textbf{both substantially outperform the original Muon recipe}.

Figure~\ref{fig:broad-context} summarizes the full six-optimizer comparison. The matched
Muon--ORBIT experiment remains the primary mechanism-isolated estimate; this broader view
shows where ORBIT sits among the frozen recipes used throughout the campaign.

\IfFileExists{figures/broad_confirmation.pdf}{%
\begin{figure}[!htbp]
\centering
\includegraphics[width=0.62\linewidth]{broad_confirmation.pdf}
\caption{Broader 124M optimizer context under independently selected recipes. Points show
mean validation loss across five seeds; error bars show one standard deviation.}
\label{fig:broad-context}
\end{figure}}{}

\FloatBarrier

\section{Discussion}
The experiments support three conclusions. First, the ORBIT improvement persists under a
strict matched-hyperparameter comparison: validation loss is lower on \MatchedWins{}
unseen seeds, and the paired confidence interval excludes zero. The cross-configuration
study complements this result by showing that ORBIT remains lower-loss under both frozen
recipes while also quantifying the influence of recipe choice on broader optimizer
comparisons.

Second, the ablations connect the empirical gain to the proposed functional geometry.
Removing the functional preconditioner produces the largest degradation, removing the RoPE
rotations produces an additional consistent degradation, and the diagonal approximation
retains most of the benefit while leaving a smaller advantage for the full $2\times2$
metric. Together these controls support the central design principle: Q/K updates benefit
from conditioning that reflects the RoPE-dependent function they induce.

Third, the improvement extends beyond the primary 124M/900-step experiment. At 2700 steps,
ORBIT attains the lowest validation loss among Muon, NorMuon, ASTRO, and ORBIT in both
tested seeds. At 355M, ORBIT remains lower than Muon and NorMuon in both tested seeds and is
competitive with ASTRO. These results position ORBIT as a practical geometry-aware
extension to Muon-family optimization whose effect is visible under matched hyperparameters
and persists across additional training regimes.

\paragraph{Scope and future work.}
The present study evaluates RoPE-modified GPT-2-style decoders with multi-head attention.
The transfer experiments use two seeds, and the 355M setting uses a smaller batch size than
the 124M setting. Extending the evaluation to grouped-query attention, QK normalization,
mixture-of-experts models, larger pretraining budgets, lower-overhead implementations, and
controlled scaling studies on substantially larger-parameter models is a direct direction
for future work. We encourage independent replication and extension of these experiments
across broader architectures and compute regimes, using the accompanying implementation
and reproducibility artifacts as a basis for further investigation of function-space
conditioning.

\section{Conclusion}
ORBIT introduces a training-time function-space preconditioner for rotary query--key
interactions. By constructing frequency-local metrics from Q/K statistics transported
through RoPE relative-position rotations, ORBIT adapts Muon's matrix update to the
computation performed by attention while leaving the model architecture and inference graph
unchanged.

Under identical hyperparameters selected from the same search grid, ORBIT lowers validation
loss on \MatchedWins{} unseen seeds with a paired mean difference of \MatchedDelta{} nats
and a 95\% confidence interval of [\MatchedCILow{}, \MatchedCIHigh{}]. Cross-configuration
experiments show that the improvement is present under both frozen recipes, and the expanded
ablations attribute the effect to functional Q/K conditioning, RoPE-aware rotation, and an additional contribution from full frequency-pair coupling. The advantage also
persists at a longer training horizon and transfers to the 355M setting against Muon and
NorMuon in the tested seeds.

These results establish ORBIT as a concrete example of computation-aware matrix
optimization: the update geometry is derived from the function of the parameter block
rather than from matrix shape alone. The same principle suggests a broader direction for
optimizers that incorporate architectural function while preserving standard model
parameters and inference-time computation.

\FloatBarrier
\bibliographystyle{plain}
\bibliography{references}
\end{document}
